\section{Vertex and bond coordinates}%
\label{sec:peps_vertex_bond_coordinates}

The virtual construction in \cite[Section~IV.C.1]{Cirac2021Matrix},
local source lines 2017--2044, can be described either by
independent labels at the vertices or by the two endpoint labels
of each edge. These descriptions are equivalent on the original
finite graph, including vertices incident to no edge.

\begin{theorem}[Endpoint-pair coordinates]%
    \label{thm:peps_vertex_bond_coordinates}
    \lean{TNLean.PEPS.VertexVirtualConfig,
        TNLean.PEPS.EdgePairVirtualConfig,
        TNLean.PEPS.vertexVirtualConfigEquivEdgePair,
        TNLean.PEPS.fullRegionVertexConfigEquiv,
        TNLean.PEPS.fullRegionVertexConfigEquivEdgePair,
        TNLean.PEPS.IsVertexVirtualConsistent,
        TNLean.PEPS.isVertexVirtualConsistent_iff_edgePair_diagonal,
        TNLean.PEPS.vertexVirtualConfigOfBonds,
        TNLean.PEPS.vertexVirtualConfigEquivEdgePair_ofBonds,
        TNLean.PEPS.isVertexVirtualConsistent_iff_exists_bonds}
    \leanok
    \uses{def:peps_virtual_bond_ground_space}
    Let every edge $e$ of a finite graph have a dimension $D_e\ge0$.
    Independent virtual configurations at all vertices are
    in bijection with configurations
    $\zeta_e=(a_e,b_e)$ of the two endpoints of each edge.
    Equality of an edge's two incident vertex labels is precisely
    the condition $a_e=b_e$.
    Such consistent configurations are exactly those obtained
    by assigning one common label to each bond.
\end{theorem}

\begin{proof}\leanok
    Read the two endpoint labels of each edge. Conversely, recover
    each vertex's incident label from its corresponding endpoint.
    An edge has distinct endpoints, so these operations are inverse.
    The consistency assertion then follows coordinate by coordinate.
\end{proof}

\begin{theorem}[The actual virtual contraction in bond coordinates]%
    \label{thm:peps_full_virtual_bond_product}
    \lean{TNLean.PEPS.fullRegionIncidentConfigEquivBonds,
        TNLean.PEPS.fullRegionVertexConfigEquiv_incident,
        TNLean.PEPS.regionVirtualBondWeight_univ_apply,
        TNLean.PEPS.fullRegionVertexVectorEquivEdgePair,
        TNLean.PEPS.fullRegionVertexVectorEquivEdgePair_virtualBondWeight}
    \leanok
    \uses{thm:peps_vertex_bond_coordinates,
        def:peps_injective_vertex_coordinates}
    In the full vertex region, sum the consistent virtual bond
    assignments in the original contraction. Under the endpoint-pair
    coordinate equivalence, the resulting coefficient vector is exactly
    \[
        \Omega(\zeta)=\prod_{e}\delta_{a_e,b_e}.
    \]
    No injectivity, connectedness, or positivity of the bond dimensions
    is assumed.
\end{theorem}

\begin{proof}\leanok
    The full region has no crossing edges, and its incident labels
    are exactly its global bond labels. An endpoint configuration
    therefore contributes once if it is consistent and zero times
    otherwise. This is the displayed product of diagonal indicators.
\end{proof}
